% Generated draft: requires mathematical and source review.
\documentclass[11pt]{article}
\usepackage[margin=25mm]{geometry}
\usepackage{fontspec}
\setmainfont{Latin Modern Roman}
\usepackage{amsmath,amssymb,mathtools,unicode-math}
\setmathfont{Latin Modern Math}
\usepackage{enumitem,needspace,xurl,hyperref}
\setlength{\parindent}{0pt}
\setlength{\parskip}{5pt}
\setlength{\emergencystretch}{3em}
\newcommand{\N}{\mathbb{N}}
\newcommand{\Z}{\mathbb{Z}}
\newcommand{\Q}{\mathbb{Q}}
\newcommand{\R}{\mathbb{R}}
\newcommand{\C}{\mathbb{C}}
\newcommand{\F}{\mathbb{F}}
\newcommand{\A}{\mathbb{A}}
\newcommand{\PP}{\mathbb{P}}
\newcommand{\E}{\mathbb{E}}
\newcommand{\eps}{\varepsilon}
\newcommand{\dd}{\,\mathrm{d}}
\begin{document}
\section*{UnsolvedMath ID: 10400203}
\textbf{Source problem code:} catalogue entry
\subsection{Definitions and mathematical statement}
For a closed oriented three-manifold \(M\), let \(\operatorname{Spin}(M)\) be its spin structures, up to isomorphism. A spin structure is a lift of the oriented frame bundle to the spin group; the set, when nonempty, is a torsor for \(H^1(M;\mathbb Z/2\mathbb Z)\) and is finite for compact \(M\).

Construct an LMO-type invariant of spin pairs, invariant under the intended spin-preserving diffeomorphisms, with the required compatibility
\[
 Z_{\mathrm{LMO}}(M)=\sum_{\sigma\in\operatorname{Spin}(M)}
 Z_{\mathrm{LMO}}(M,\sigma).
\]
The excerpt does not prescribe its value space, normalization, surgery construction or additional universal/finite-type properties. The sum identity alone does not determine a unique refinement, so no arbitrary splitting is asserted to fulfill the intended natural construction. The source's closed-manifold scope is not restricted to rational homology spheres.

The source also requests a similar refinement for \(\operatorname{spin}^c\) structures but gives no explicit formula. Those structures form a torsor for \(H^2(M;\mathbb Z)\), which may be infinite. Thus an analogous sum over all such structures would need finite support, a specified completion or another summability convention; none is silently supplied.

The registry has an open label, while the summary calls this partial progress and says spin LMO invariants were defined by various authors without naming constructions or references. These source reports do not verify a construction with the exact displayed compatibility for all intended manifolds.
\subsection{Short English statement}
Refine LMO by a spin structure so summing the refinements recovers the original invariant, and formulate the analogous \(\operatorname{spin}^c\) theory. The value space and naturality requirements are missing, and an infinite \(\operatorname{spin}^c\) sum needs a separate convention.
\subsection{Sources}
\begin{itemize}
\item Local registry: \url{data/unsolvedmath/problems.json}, permanent UnsolvedMath ID \(10400203\), source code catalogue entry; supplied statement, background, and research summary, updated 2026-08-06.
\item Ohtsuki, \emph{Problems on invariants of knots and 3-manifolds} (2002), Problem 11.7, attributed to T. Le and V. Turaev, PDF page 156: \url{https://msp.org/gtm/2002/04/gtm-2002-04-024s.pdf}. Supplied PDF-font-structure extraction accessed 2026-07-29.
\item The summary reports spin-LMO constructions by unnamed authors, without exact publications or theorem statements.
\end{itemize}
\end{document}
